\section{Parallel treatments in Aquinas}
\label{app:parallels}

The following loci develop the distinctions used in the systematic exposition.
The \emph{Summa}'s connected account supplies the sequence; the other works
make particular arguments more explicit.

\begin{longtable}{@{}>{\raggedright\arraybackslash}p{0.28\linewidth}>{\raggedright\arraybackslash}p{0.24\linewidth}>{\raggedright\arraybackslash}p{0.41\linewidth}@{}}
\toprule
Work and locus & Related \emph{Summa} loci & Contribution\\
\midrule\endhead
\emph{De ente et essentia}, the discussion beginning \emph{Nunc restat
videre} and its sequel \emph{His igitur visis} & I q.~50 a.~2 ad~3 &
Separate substances lack matter while retaining the reception of existence.
The Baur--Koch/Corpus Thomisticum witness numbers these chapters 3--4;
other editions divide the work differently.\\
\emph{De veritate}, q.~8 & I qq.~54--58 & Angelic knowledge: its modes,
objects and limits; the relation of natural intellect to the vision of God.\\
\emph{De veritate}, q.~9 aa.~1--4 & I qq.~106--107 & Illumination,
purification and perfection as the communication and reception of knowledge;
speech as voluntary disclosure.\\
\emph{De malo}, q.~16 & I qq.~63--64, 109, 114 & The demons' nature,
sin, persistence in evil, knowledge and activity. The question develops the
relation between enduring natural goods and the evil of the will.\\
\emph{Summa}, III q.~8 a.~4 & I q.~108 a.~8 & Christ as head of angels
and human beings within one mystical body, with distinct ways of receiving
his gifts.\\
\emph{Summa}, III q.~30 aa.~1--4 & I q.~51 a.~2;
q.~108 a.~5; q.~112 a.~1 & The Annunciation: the messenger, Mary's
consent, Gabriel's commission, the bodily appearance and the order of the message.\\
\bottomrule
\end{longtable}

The parallel texts clarify the same underlying relation: the creature's
highest perfection is received. The act of being, intelligible truth,
supernatural beatitude and a ministry to another each require their proper
account; none is explained merely by calling the recipient immaterial.

\subsection{The shared light in \emph{De veritate}}

The earlier disputation makes the recipient's activity especially clear.
A teacher offers a principle through which the pupil can understand something
previously beyond his grasp. Angelic illumination accomplishes this without
vocal sounds or successive reasoning: one intellect communicates truth
according to another's capacity. The lower angel receives a real perfection
and truly understands. The higher angel neither creates that intellect nor
bestows its beatitude.\footnote{\emph{De veritate}, q.~9 a.~1,
corpus and ad~2--4, 9--10, 13, 16; Latin Leonine-aligned electronic witness,
Corpus Thomisticum.}

Purification, illumination and perfection name distinguishable moments in
the reception of intelligible truth. Purification removes ignorance;
illumination supplies the means of knowing; perfection is the knowledge
attained. A blessed angel can thus receive a new particular revelation while
already enjoying God. Holiness is compatible with receptivity to further
gifts.\footnote{\emph{De veritate}, q.~9 a.~3, corpus and ad~1--5.}

Speech expresses the will to share. An angel's nature can be known without
every actual thought becoming public. The speaker directs a conception to
another, distinguishing possession of an intelligible form, actual
consideration, and communication. This willing disclosure gives spiritual
communion a personal form: understanding is shared by distinct knowers who
address one another.\footnote{\emph{De veritate}, q.~9 a.~4, corpus
and ad~7--9, 11, 15.}
